\documentclass[10pt,a4paper]{amsart}
\usepackage[margin=19mm]{geometry}
\usepackage{amsmath,amssymb,mathtools}
\usepackage[scr=rsfs,cal=euler]{mathalfa}
\usepackage{xfrac}
\usepackage[unicode,hidelinks,bookmarks=false]{hyperref}
\newcommand{\ie}{i.e.\ }
\newcommand{\cf}{cf.\ }
\newcommand{\resp}{respectively\ }
\newcommand{\Subsection}{\subsection}
\providecommand{\nobreakdash}{\nobreak\hbox{-}}
\setlength{\emergencystretch}{2em}
\multlinegap0pt
\usepackage{bm}
\usepackage{bbm}
\usepackage{dsfont}
\usepackage{xspace}
\usepackage{cancel}
\usepackage{mathtools}
\usepackage{amsthm}
\makeatletter
\@ifundefined{theorem}{\newtheorem{theorem}{Theorem}}{}
\@ifundefined{proposition}{\newtheorem{proposition}[theorem]{Proposition}}{}
\makeatother
\title[Real gamma distribution on analytic bundles of flag varietie]{Real gamma distribution on analytic bundles of flag varieties}
\author{Haoming Wang}
\date{Source: arXiv:2601.21304v2 (CC BY 4.0)}
\begin{document}
\maketitle

\noindent\emph{Source and licence.} Real gamma distribution on analytic bundles of flag varieties. Version arXiv:2601.21304v2 (CC BY 4.0), distributed under the Creative Commons Attribution 4.0 International licence, as recorded in the arXiv licence metadata for this exact version and shown on the abstract page https://arxiv.org/abs/2601.21304v2. Full rights notice: licenses/arxiv-2601.21304-CC-BY-4.0.txt.\par\smallskip

\noindent\textit{Editorial scope.} Counted unit: the Theorem whose source label is \texttt{thm: main theorem} in the article (source0.tex lines 561--571, source label \texttt{thm: main theorem}), delivered together with the article's own complete original proof of it (source0.tex lines 594--621, one \texttt{proof} environment of 28 lines). No mathematics is written, repaired or re-derived: statement and proof are the article's own text line for line. The proof is detached from the statement in the source: the article states the result at lines 561--571 and prints its proof at lines 594--621, and the proof environment's own header names the statement, so the association is the article's own. Required context retained uncounted: lem: hypergeometric (lines 278--284); lem: james55 (lines 287--293); lem: Khatri66 (lines 297--299); prop: herz55 (lines 329--331). Every definition, display and label of those lines is retained unchanged. Omitted from this package and not redistributed: 1--277, 285--286, 294--296, 300--328, 332--560, 572--593, 622--880. Nothing inside a retained excerpt is omitted or altered: every retained body is delivered exactly as the article's lines are written, so no omission receipt of any kind accompanies this package. Citations: the retained excerpts carry no citation command at all, so no bibliography entry of the article is transcribed and no reference is redistributed. Reference adaptation: none was needed and none was made; every reference inside the delivered unit resolves inside it, so no unresolved reference or citation is produced. Printed slips kept literal: no printed word, symbol, hypothesis or number of the counted statement or of its proof was altered. Layout and class: the article's own class is \texttt{amsart} with packages including graphicx, multirow, amsmath, amssymb, amsfonts, amsthm, mathrsfs, appendix, xcolor, textcomp, manyfoot, booktabs, algorithm, algorithmicx; the wrapper is the lane's pinned \texttt{amsart} editorial wrapper at 10pt on A4, which restates the theorem-like environments the retained text needs and adds the article's own macro definitions that the retained text needs (none), with extra wrapper packages (bm, bbm, dsfont, xspace, cancel, mathtools). The delivered numbering is the unit's own. Assets: no retained span contains a figure environment, a graphics-inclusion command, a PDF-page inclusion, a table or a TikZ picture; no image file is copied into this package, the package declares no asset dependency and no image file is redistributed with it. Licence: version arXiv:2601.21304v2 of this article is distributed under the Creative Commons Attribution 4.0 International licence, as recorded in the arXiv licence metadata for this exact version and shown on the abstract page https://arxiv.org/abs/2601.21304v2. A full rights notice is delivered at licenses/arxiv-2601.21304-CC-BY-4.0.txt. Characters: every retained excerpt is pure ASCII, so the delivered engine, XeLaTeX with its default Latin Modern text font, reports no missing character.
\begin{proposition}\label{lem: hypergeometric} Suppose $ A$ is an $n\times n$ and $ B$ is a $k\times k$ symmetric matrix. 
    \[\begin{aligned}
        \int_{O_n}{}_{0}F_{0}( A H_1 B H_1^{\prime}) [dH] = {}_{0}F_{0}( A, B),\\ 
         H = (H_{1},  H_{2}), \text{ where }  H_1 \text{ is } n \times k,
    \end{aligned}\]
    that is, the integral runs over the orthogonal group $O_n$.
\end{proposition}

\begin{proposition}\label{lem: james55} Suppose $X$ is an $n\times k$ matrix.
    \[\begin{aligned}
        \int_{O_n} \operatorname{etr}( X H_1^{\prime}) (d H) = {}_{0}F_{1}\left(\frac{1}{2}n;\frac{1}{4} X' X\right), \\ 
         H = (H_{1},  H_{2}), \text{ where }  H_1 \text{ is } n \times k,
        \end{aligned}\]
    and the integral runs over the orthogonal group $O_n$ similar to Lemma \ref{lem: hypergeometric}.
\end{proposition}

		\begin{equation*}
			\int_{ X^{\prime}  X= S} \operatorname{etr} ( A  X  B  X^{\prime}) (d  X)= \frac{\pi^{\frac{nk}{2}}}{\varGamma_{p}(\frac{n}{2} )} | S|^{\frac{n-k-1}{2}} {_0}F_0( A,  B S).\label{eq: Khatri1966}
		\end{equation*}\label{lem: Khatri66}

\begin{proposition}\label{prop: herz55}
$(dZ) = 2^{-k} |R|^{\frac{1}{2}(n-k-1)} \cdot (dR) \cdot (dK)$, where $(dK) = {(H'dH)} \cdot {(H_\perp'dH)}$.
\end{proposition}

\begin{theorem} \label{thm: main theorem}  Suppose $X$ is an $n\times k$  real normal matrix $T_1$ with $B = ( b_{1}, b_{2},\dots, b_k)$,
%\[ \varSigma_{1} = \sum_{i=1}^{p} \sum_{j=1}^{p} A_{ij} \otimes  B_{ij},\quad   A_{ij} = ( I_{n}\otimes  b_{i})\Sigma_{1}( I_{n}\otimes  b_{i}^{\prime}),\quad  B_{ij} = \overline{ b_{i}}^{\prime} b_{j}.\]
Let $ M$ be a fixed real $n\times k$ matrix and $ S =  (X +  M)'(X +  M)$. Assume that the diagonals in the $k\times k$ real matrix $ U = (u_{ij})$ with $u_{ij} = \operatorname{tr}( A_{ij})$ are all positive, that is, $u_{ii}>0 \, (i=1,2,\dots,k)$. Then the probability density distribution of $ S$ only depends on $ T %=  B^{\prime} S B 
= (t_{ij})$ with 
    $t_{ij} = \operatorname{tr}( B_{ij}  S)$ when $n > k - 1$, and is 
        \begin{eqnarray}
         \frac{| \varTheta_{1}|^{\frac{1}{2}}}{2^{\frac{nk}{2}}\varGamma_{k}(\frac{n}{2})} \operatorname{etr} \left(-\frac{1}{2}  U T \right) | T|^{\frac{n-k-1}{2}}          \text{ when } {M}=0; \\
                \times \operatorname{etr} \left(-\frac{1}{2}  \varOmega\right){}_{0}F_{1}\left(\frac{n}{2};\frac{1}{4} \varDelta  T\right) \text{ when } {M}\neq0;
        \end{eqnarray}
        where $\varOmega = \sum_{i,j=1}^{k} B_{ij}  M^{\prime}   A_{ij}  M$ and $ \varDelta = \sum_{i,j,k,l=1}^{k} B'_{ij} M'A'_{ij}A_{kl} M B_{kl}$.
\end{theorem}

\begin{proof}[Proof of Theorem \ref{thm: main theorem}]
    Suppose the central part holds for $n > k-1$. Decomposing $ X =  H  Z$ where $ H' H =  I_p$.
    Let $K = (H,H_{\perp}) \in O(n)$. From Lemma \ref{prop: herz55}, we could rewrite the density in $T_1$ as 
    \[\begin{aligned}
        \operatorname{etr} \left(-\frac{1}{2}  \varOmega\right)\int_{O_n}\operatorname{etr}\left(-\frac{1}{2} \sum_{i,j=1}^{k}( A_{ij} M B_{ij})Z'H' \right) (d K) \\
        = \operatorname{etr} \left(-\frac{1}{2}  \varOmega\right){}_{0}F_{1}\left(\frac{n}{2};\frac{1}{4}\sum_{i,j,k,l=1}^{k}( B'_{ij} M'A'_{ij}A_{kl} M B_{kl}) Z'Z\right).    \end{aligned}\]
    %where we use the fact that these $ B_{ij}$ are orthogonal to each other,\begin{equation}   \begin{aligned}         B_{ij}  B_{kl} =   B_{il} \delta_{jk}, \quad   B_{ij}^{\prime} =   B_{ji}, \quad   A_{ij}^{\prime} =   A_{ji}.    \end{aligned}    \label{eq: commutativity}    \end{equation}
    Thus, if the central part is true, Theorem \ref{thm: main theorem} is consequently proved by applying Lemma \ref{lem: james55}. 
    
    From Lemma \ref{lem: Khatri66} on the quadratic forms in normal vectors, we can reduce the central part of Theorem \ref{thm: main theorem} to the $\frac{1}{2}k(k+1)$ independent quadratic forms by introducing $ Y =  X B = ( y_{1}, y_{2}\dots,  y_{k})$
    \begin{equation*}
    \begin{aligned}
        {y}_1'{A}_{11}{y}_1, {y}_1'{A}_{12}{y}_2, \dots, {y}_1'{A}_{1k}{y}_k,\\
        {y}_2'{A}_{22}{y}_2, \dots, {y}_2'{A}_{2k}{y}_k, \\
        \vdots\qquad \\
        {y}_k'{A}_{kk}{y}_k,
    \end{aligned}
\end{equation*}
In fact, we have for each ${y}_i'{A}_{ij}{y}_j$ the contribution to the probability density function
\[\begin{aligned}
    {_{0}}F_{0}\left( I - \frac{q_{ij}}{2} A_{ij}, q_{ij}^{-1} y_{i} y_{j}^{\prime}\right) = {_{0}}F_{0}\left( I - \frac{q_{ij}}{2} A_{ij}, q_{ij}^{-1} y_{j}^{\prime} y_{i}\right).\end{aligned}\]
However, we may find $ y_{j}^{\prime} y_{i} = t_{ij} = \operatorname{tr} ( B_{ij} S)$ so that from these properties of hypergeometric functions
\[\begin{aligned}
    {}_{0}F_{0} ( X,c Y) = {}_{0}F_{0} (c X, Y)\\
    {}_{0}F_{0} ( X,  I) = {}_{0}F_{0} ( X) = \operatorname{etr}( X), 
\end{aligned}\]
the terms concerning $q_{ij}$ cancel out. This yields the desired form.
\end{proof}

\end{document}
